% --- Patron: Ikasha, She Who Sleeps (Latent Potential & Shadow) ---
\subsection{Ikasha, She Who Sleeps (Latent Potential \& Shadow)}

\textit{Lore.} Ikasha is the hush between footfalls, the patience of dark water, the black‑feathered watcher at every threshold. In stillness she gathers what might be, at crossroads she whispers of what may yet come. Ravens circle her, bearing secrets between worlds. Her followers learn to move unseen and speak unremembered, becoming shadows that slip between what is and what could be.

Ikasha does not act; she waits. She does not command; she suggests. Her power is not in force but in possibility – the hidden door, the forgotten name, the secret that, if spoken, would change everything. Those who serve her become keepers of thresholds, listeners at the edge of sleep, and witnesses to the shape of things not yet born. She is the patron of spies, dreamers, outcasts, and those who know that the darkest shadow is also the safest cloak.

\begin{quote}
“Blow out the candle. If the room listens back, ask softly. At the next crossroads, the raven waits – and the shadow remembers your passing.”
\end{quote}

\noindent\textbf{Domain Focus}
\begin{itemize}
  \item \textbf{Latent Potential:} the power of what is not yet, the seed of action, the whisper before the scream
  \item \textbf{Shadow as Shelter:} concealment, secrets, and the spaces between notice
  \item \textbf{Thresholds of Change:} doors, crossroads, dreams, and the boundaries of sleep
  \item \textbf{Raven’s Witness:} omens, messages, and the weight of secrets shared
\end{itemize}

\subsubsection*{Rites of Ikasha}

\paragraph{Rite of the Unlit Candle (Low, 4 XP)} \hfill \\
\emph{Action; Self; Yes (Stealth).} \\
\textbf{Materials:} A candle snuffed with a whisper. \\
\textbf{Effect:} You step into the space between notices. Gain \textbf{+1 die} to a single \textit{Stealth} or \textit{Subterfuge} roll, and if you succeed, you may also declare one minor detail “forgotten” by any observer (e.g., “they don’t remember my face,” “they left the key in the lock”). \\
\textbf{Push It:} Extend the effect to one ally within Near, but your shadow lingers independently for a moment – the GM gains \textbf{1 SB (Diamonds)} as evidence remains. \\
\emph{Requires: Familiar (\textit{Invoke:} 1 Boon).}

\paragraph{Rite of the Crossroads Raven (Low, 5 XP)} \hfill \\
\emph{Scene; Near; No.} \\
\textbf{Materials:} Three black feathers scattered at a crossing. \\
\textbf{Effect:} A raven appears, visible only to you and your allies. It offers an omen: you may reroll one failed navigation, pursuit, or escape roll this scene, or force one enemy to hesitate (lose their next minor action). \\
\textbf{Push It:} The raven whispers a true secret about the current situation (the GM must answer one yes/no question truthfully). In return, it demands a secret from you – the GM gains \textbf{1 SB (Hearts)} to spend later as a complication involving that secret. \\
\emph{Requires: Familiar (\textit{Invoke:} 1 Boon).}

\paragraph{Rite of the Umbral Reservoir (Standard, 8 XP)} \hfill \\
\emph{Action; Self or Ally; No.} \\
\textbf{Materials:} Water gathered on a moonless night, held in a black vessel. \\
\textbf{Effect:} Draw on deep reserves of patience and shadow. Choose one:
\begin{itemize}
  \item Gain \textbf{+2 dice} on your next \textit{Stealth}, \textit{Deception}, or \textit{Resolve} roll.
  \item Clear \textbf{1 Fatigue} from yourself or a touched ally.
  \item Gain one free success on an attempt to escape restraint or pursuit (no roll needed).
\end{itemize}
\textbf{Push It:} Apply two benefits instead of one, but you mark \textbf{+1 Obligation} and the GM gains \textbf{1 SB (Clubs)} as the shadow’s debt comes due. \\
\emph{Requires: Familiar + Codex (\textit{Invoke:} 1 Boon).}

\paragraph{Rite of the Secret Keeper (Standard, 9 XP)} \hfill \\
\emph{Instant; Touch; No.} \\
\textbf{Materials:} A lock of hair, a personal token, or a name written on black paper. \\
\textbf{Effect:} You compel a truthful answer to one direct question about a secret the target knows. They may resist with a Resolve test (DV 4); if they fail, they must answer truthfully (though they may be cryptic). If they succeed, they suffer \textbf{1 Fatigue} from the effort of concealment. \\
\textbf{Push It:} You also learn a hidden emotion or related secret (GM’s choice). In exchange, the target learns one of your secrets – the GM gains \textbf{1 SB (Hearts)} and may reveal that secret at a dramatically appropriate moment. \\
\emph{Requires: Familiar + Codex (\textit{Invoke:} 1 Boon).}

\paragraph{Rite of the Unlatched Door (Standard, 7 XP)} \hfill \\
\emph{Instant; Near; Yes.} \\
\textbf{Materials:} A key that has never been used, bent in a circle. \\
\textbf{Effect:} You persuade a threshold to open. One mundane lock, seal, or barrier (door, chain, shutter, even a pursuer’s grip) opens or loosens as if it were never secured. This does not work on magical wards. \\
\textbf{Push It:} You may also silence any noise the opening would make, but the threshold “remembers” you – the next time you pass this way, the GM gains \textbf{1 SB (Spades)}. \\
\emph{Requires: Familiar + Codex (\textit{Invoke:} 1 Boon).}

\paragraph{Rite of the Shadow at the Crossroads (High, 12 XP)} \hfill \\
\emph{Scene; Self; No.} \\
\textbf{Materials:} Stand in absolute darkness or at a deserted crossroads at midnight. \\
\textbf{Effect:} You become a living shadow. For the scene, you are intangible to mundane harm, can pass through small gaps and thresholds, gain \textbf{+2 dice} to all \textit{Stealth} rolls, and automatically succeed on one attempt to escape restraint or pursuit. You cannot manipulate objects normally (though you can open unlocked doors). \\
\textbf{Push It:} You may interact once with a locked or warded object (pick a lock, lift a seal, read a sealed letter) while intangible, but you become partially corporeal for a beat – your next defence roll is at \textit{Desperate} position. Mark \textbf{1 SB (Spades)} as reality protests. \\
\emph{Requires: Familiar + Codex + Tier III (\textit{Invoke:} \textbf{2 Boons}).} \\
\emph{Obligation:} 7 segments.

\paragraph{Rite of the Unfinished Dream (High, 14 XP)} \hfill \\
\emph{Extended; Self; No.} \\
\textbf{Materials:} A dream you remember vividly, written down and burned. \\
\textbf{Effect:} You reach into the realm of latent potential. Choose one:
\begin{itemize}
  \item \textbf{Whisper of Possibility:} Once before the next dawn, you may reroll any single roll after seeing the result (but before the outcome is declared).
  \item \textbf{Hidden Path:} You and your allies may bypass one major obstacle (a guarded gate, a collapsed bridge, a social dead end) as if a secret way had always been there – but the GM gains \textbf{1 SB} as the world adjusts.
  \item \textbf{Unspoken Name:} You learn the true name of one creature or place (GM discretion). That name can be used for future rites or bargains.
\end{itemize}
\textbf{Push It:} Gain two of the above benefits, but you mark \textbf{+2 Obligation} and the GM gains \textbf{2 SB} (the dream’s echoes attract attention). \\
\emph{Requires: Familiar + Codex + Tier III (\textit{Invoke:} \textbf{2 Boons}).} \\
\emph{Obligation:} 7 segments.

\subsubsection*{Cantors \& Cults of Ikasha}

\textbf{The Cantor’s Song.} A Cantor of Ikasha sings in \textit{whispers} – melodies that are felt rather than heard, harmonies that slip beneath the threshold of attention. Their voice is used to hide a crowd’s escape, to make a guard forget why they entered a room, or to pass a message through enemy lines without a single word being spoken aloud. Ikasha’s Cantors are rarely recognised; they are the ones who hum absently while others plan rebellion.

\textbf{The Cult of the Unlit Candle.} Ikasha’s cults are shadowy by necessity, meeting in basements, abandoned bell‑towers, and fog‑bound crossroads. They whisper freedom to the oppressed – escaped slaves, hunted heretics, debtors fleeing their creditors. A ritual of the Unlit Candle involves each member extinguishing a flame and speaking a secret wish for escape. The candles are never relit; the wish is carried away on ravens’ wings. Their creed is never written, only passed from shadow to shadow: \textit{“The dark does not confine – it conceals the road that the light would forbid.”} Outsiders who stumble into a meeting are offered a black feather and a choice: forget, or follow.

\subsubsection*{Witchcraft of Ikasha (Threshold / Shadow)}

A witch who walks with Ikasha is a \textit{threshold‑keeper} or \textit{shadow‑weaver} – one who knows that a door left unlatched is not an invitation to theft, but to liberation. Their magic works in the gaps: the crack under a door, the moment between heartbeats, the pause between a guard’s patrol. This craft can be invisible – a bent lock, a misplaced key, a shadow that moves of its own accord – or as grand as erasing a name from a pursuer’s memory.

\textbf{The Witch’s Tool: The Raven’s Quill.} A quill carved from a raven’s primary feather, dipped not in ink but in the absence of light. Once per session, when used to write a name on a threshold (a doorframe, a gatepost, a windowsill), that name becomes “forgotten” for a scene: the person whose name is written cannot be perceived by those who cross that threshold until they succeed on a Resolve test (DV 3). The quill can also erase a single word from a written document (it becomes a smudge of shadow).

\textbf{A Hedge Gift: Unlatched (2 XP).} Once per scene, you may declare a mundane lock, latch, or binding to be “unlatched” – it opens freely (or fails to close) as if it had never been secured. This works on doors, chains, shutters, and even a pursuer’s grip, but not on magical wards.

\subsubsection*{Ikasha’s Corruption Table}

\begin{tblr}{
  colspec = {Q[c,1.5cm] X[2.5] X[2.5]},
  rowhead = 1,
  row{odd} = {bg=gray!10},
  row{1} = {bg=gray!30, font=\bfseries},
  hlines,
}
Tier & Benefit & Cost / Quirk \\
1 & \textbf{Shadow’s Whisper:} +1 die to \textit{Deception} when lying in darkness or through intermediaries. & \textbf{Truth Echoes:} Occasionally speak in riddles or half‑truths without realising it. \\
2 & \textbf{Umbral Sight:} Once per scene, gain +2 dice to \textit{Notice} hidden threats or ambushes in dim light. & \textbf{Light Sensitivity:} –1 die to rolls requiring keen vision in bright conditions. \\
3 & \textbf{Raven’s Memory:} Never forget a secret told to you; recall any whispered conversation perfectly. & \textbf{Secret Burden:} Must keep one troubling secret per Tier; suffer \textbf{1 Fatigue} when actively trying to forget. \\
4 & \textbf{Threshold Walker:} Once per session, move through one locked door, sealed letter, or magical barrier without triggering wards. & \textbf{Crossroads Bound:} Must pause at literal or metaphorical crossroads to “listen” before major decisions. \\
5 & \textbf{Shadow Debt:} Once per session, call in a favour from someone who owes you a secret – no questions asked. & \textbf{Obligation Web:} Every secret you learn creates a subtle tie to its keeper; mark \textbf{1 SB (Diamonds)} when ignoring these connections. \\
6+ & \textbf{Umbral Ascendancy:} Once per session, become completely undetectable to all non‑magical senses for one scene. & \textbf{Reality Thinning:} The boundary between what is and what could be grows thin – the GM may introduce subtle reality alterations around you. \\
\end{tblr}

\subsection*{Playstyle Notes}
Ikasha rewards patience, subtlety, and the willingness to let things unfold. Her followers are not agents of action but of possibility – they open doors, plant secrets, and wait for the harvest. Intrusions often take the form of dreams showing a hidden danger, a raven appearing with a cryptic message, or a secret you thought buried resurfacing at the worst moment. Refusing an intrusion does not bring wrath – merely the withdrawal of Ikasha’s favour: doors stay locked, shadows become just shadows, and the road ahead grows a little darker.

\noindent\textbf{Emphasises}
\begin{itemize}
  \item \textbf{Latent Potential:} power as possibility, not force
  \item \textbf{Shadow as Shelter:} concealment, secrets, and the spaces between notice
  \item \textbf{Thresholds of Change:} doors, crossroads, dreams, and the boundaries of sleep
  \item \textbf{Raven’s Witness:} omens, messages, and the weight of secrets shared
\end{itemize}